% !TEX program = xelatex
% XeLaTeX；A4 黑白双栏。正文、习题、答案自然分页；各部分从奇数页开始。
\documentclass[UTF8,10pt,a4paper,twocolumn,twoside,fontset=fandol]{ctexart}
\usepackage[a4paper,left=15mm,right=15mm,top=16mm,bottom=16mm,columnsep=8mm,headheight=14pt,headsep=6mm,footskip=9mm]{geometry}
\usepackage{amsmath,amssymb,xcolor,fancyhdr,tikz,needspace}
\ifdefined\XeTeXversion\else
  \errmessage{Please compile this document with XeLaTeX}
\fi
\pagestyle{fancy}
\fancyhf{}
\fancyfoot[C]{\footnotesize\thepage}
\renewcommand{\headrulewidth}{0pt}
\renewcommand{\footrulewidth}{0pt}
\setlength{\parindent}{1.6em}
\setlength{\parskip}{2pt}
\setlength{\abovedisplayskip}{5pt}
\setlength{\belowdisplayskip}{5pt}
\setlength{\abovedisplayshortskip}{3pt}
\setlength{\belowdisplayshortskip}{3pt}
\setlength{\fboxsep}{5pt}
\setlength{\fboxrule}{0.35pt}
\raggedbottom
\newcommand{\R}{\mathbb R}
\newcommand{\eps}{\varepsilon}
\newcommand{\limn}{\lim_{n\to\infty}}
\newcommand{\textspacing}{%
\fontsize{10}{15}\selectfont
\setlength{\parskip}{1.5pt}% minipage 会重设段距，须在栏内指定。
\setlength{\abovedisplayskip}{4pt}
\setlength{\belowdisplayskip}{4pt}
\setlength{\abovedisplayshortskip}{3pt}
\setlength{\belowdisplayshortskip}{3pt}}
\newcommand{\head}[1]{\par\addvspace{5pt}\noindent{\fontsize{13}{18}\selectfont\bfseries #1}\par\nobreak\vspace{3pt}}
\newcommand{\subhead}[1]{\par\addvspace{6pt}\noindent{\fontsize{11.5}{16}\selectfont\bfseries #1}\par\nobreak\vspace{1pt}}
\def\exparts#1\quad#2\endexparts{{\fontsize{10.5}{15}\selectfont\bfseries #1}\hspace{0.8em}{\normalfont\fontsize{10}{15}\selectfont #2}}
\newcommand{\ex}[1]{\par\addvspace{5pt}\noindent\exparts#1\endexparts\par\nobreak}
\newcommand{\prooflead}[1]{\par\addvspace{3pt}\noindent{\fontsize{9}{14}\selectfont\bfseries #1}\hspace{0.7em}\ignorespaces}
\newcommand{\badge}[1]{\tikz[baseline=(tag.base)]\node[rounded corners=2pt,draw=black!65,fill=white,line width=0.4pt,inner xsep=3pt,inner ysep=1pt,font=\fontsize{9}{11}\selectfont\bfseries,text=black] (tag) {#1};}
% 模块标题：类别为黑色粗体，名称用常规字重；固定间隔代替手敲空格。
\newcommand{\modulehead}[2]{\badge{#1}\hspace{0.7em}{\normalfont\fontsize{10}{15}\selectfont\color{black}#2}}
\newcommand{\lead}[1]{\par\addvspace{3pt}\noindent\badge{#1}\hspace{0.7em}\ignorespaces}
\newcommand{\graybox}[2]{\par\addvspace{5pt}\noindent
\begingroup\setlength{\fboxsep}{0pt}% 标题条贴合外框，正文单独设置内边距。
\fcolorbox{black!65}{black!2}{%
\begin{minipage}{\dimexpr\linewidth-2\fboxrule\relax}
\textspacing
\noindent{\setlength{\fboxsep}{3pt}%
\colorbox{black!12}{\parbox{\dimexpr\linewidth-6pt\relax}{\hspace{2pt}\color{black}\bfseries #1}}}\par\nointerlineskip
\vspace{5pt}\noindent\hspace*{5pt}%
\begin{minipage}{\dimexpr\linewidth-10pt\relax}
\textspacing #2\par
\end{minipage}\par\vspace{5pt}
\end{minipage}}\endgroup\par\addvspace{5pt}}
\newcommand{\q}[2]{\par\addvspace{7pt}\noindent\textbf{#1.}\enspace #2\par}
\newcommand{\work}[1]{\vspace{0pt plus #1fill}}
\makeatletter
% 正文采用自然分页，不用强制换页。
\newcommand{\evenbreak}{%
  \clearpage
  \ifodd\value{page}\else
    \hbox{}\thispagestyle{empty}\clearpage
  \fi}
\newcommand{\evenfinish}{%
  \clearpage
  \ifodd\value{page}\hbox{}\thispagestyle{empty}\clearpage\fi}
\NewDocumentCommand{\textcol}{m +m}{%
\textspacing #2\par}
\NewDocumentCommand{\quizcol}{m m +m}{%
\noindent
\begin{minipage}[t][\dimexpr\@colht-12pt\relax][s]{\columnwidth}
\fontsize{10}{15}\selectfont
\graybox{#2}{\small 建议20分钟。写出主要步骤；空白处可作草稿。}
#3
\par\noindent{\footnotesize 用时：\underline{\hspace{12mm}}\quad 待订正题号：\underline{\hspace{16mm}}}
\end{minipage}\par}
\makeatother
\begin{document}
\twocolumn[{\centering\fontsize{14}{20}\selectfont\bfseries
数学分析小测复习讲义\par\vspace{2pt}
{\normalfont\fontsize{9}{12}\selectfont 中国科学技术大学\quad 缪语博\quad PB26000200\par}\vspace{4pt}}]

% 第1页：概念及定义证明。
\textcol{1}{
\head{1\quad 实数与极限}
{\footnotesize 定义、定理编号沿用课本；框外另作解释与应用。\par}
\subhead{确界与实数的完备性}
\graybox{\modulehead{定义 1.10}{上确界与下确界}}{
设 $a$ 是数集 $X$ 的上界，若对于任意的正数 $\eps$，都存在一个数 $x_\eps\in X$，使得 $x_\eps>a-\eps$。则称 $a$ 为 $X$ 的上确界，记为 $\sup X$。同理可定义下确界，记为 $\inf X$。}
\prooflead{理解}上确界是最小的上界。证明时先说明它是上界，再说明任何更小的数都挡不住整个集合。上确界属于集合时，才是最大值。
\ex{例1\quad 求确界并说明理由}
设 $E=\{1-1/n:n=1,2,\ldots\}$。每个元素都小于 $1$，所以 $1$ 是上界。任给 $\eps>0$，取正整数 $n>1/\eps$，则
\[
 1-\frac1n>1-\eps.
\]
因此 $\sup E=1$，但 $E$ 没有最大值。又 $0\in E$ 且所有元素非负，故 $\inf E=\min E=0$。
\subhead{数列极限的定义}
\graybox{\modulehead{定义 1.2}{数列极限}}{
设 $\{a_n\}$ 是给定的数列。如果有一个实数 $a$ 具有下列性质：对任意给定的一个正数 $\eps$，总是存在一个自然数 $N=N(\eps)$，使得当 $n>N$ 时，不等式 $|a_n-a|<\eps$ 即 $a-\eps<a_n<a+\eps$ 成立，那么称 $a$ 是数列 $\{a_n\}$ 的极限，或数列 $\{a_n\}$ 以 $a$ 为极限，也称数列 $\{a_n\}$ 收敛于 $a$，记成
\[
\lim_{n\to\infty}a_n=a\text{ 或 }\lim a_n=a,
\]
也可以记成 $a_n\to a\ (n\to\infty)$，读作“当 $n$ 趋于无穷时，$a_n$ 趋于 $a$”。}
\prooflead{理解}$|a_n-a|$ 是两点的距离。先提出误差要求 $\eps$，再选 $N$；之后的\textbf{每一项}都要满足要求。$N$ 不必最小。课本用 $a$ 表示极限，例题中的 $A$ 只是另一个字母。
\graybox{易错点}{
“有无穷多项接近 $A$”不等于收敛到 $A$。如 $0,1,0,1,\ldots$ 虽有无穷多项等于 $0$，但后面始终还有等于 $1$ 的项。}
}
% 自然分页：不人为插入分页。
\textcol{1}{
\subhead{定义证明：先估计误差，再选项号}
\ex{例2\quad 证明 $\displaystyle\frac{3n+1}{n+2}\to3$}
\prooflead{分析}先算误差：
\[
 \left|\frac{3n+1}{n+2}-3\right|
 =\frac5{n+2}<\frac5n.
\]
只要 $n>5/\eps$ 就够了。寻找 $N$ 时可以放大误差，换来一个容易控制的式子。
\prooflead{证明}任给 $\eps>0$，取正整数 $N>5/\eps$。当 $n>N$ 时，
\[
 \left|\frac{3n+1}{n+2}-3\right|<\frac5n<\eps.
\]
由定义，所求极限为 $3$。
\subhead{收敛数列的基本性质}
\graybox{\modulehead{定理 1.3}{极限的基本性质}}{
$1^\circ$ 若 $\{a_n\}$ 是收敛数列，则 $\{a_n\}$ 的极限是唯一的。
\par $2^\circ$ 改变数列中有限多项的值，不影响数列的敛散性及其极限。}
\lead{定理 1.8}设 $\{a_n\}$ 收敛于 $a$，则其任意一个子数列也收敛于 $a$。
\prooflead{定理 1.4（节选）}设 $\{a_n\}$ 为收敛数列，极限值记为 $a$。则 $1^\circ$ $\{a_n\}$ 为有界数列。即存在一个正数 $M$（与变量 $n$ 无关），使得 $|a_n|\le M$ 对所有的 $n=1,2,3,\cdots$ 成立。
\prooflead{证明}设 $a_n\to A$，在定义中取 $\eps=1$，可得某个 $N$，使 $n>N$ 时 $|a_n|<|A|+1$。令
\[
 M=\max\{|a_1|,\ldots,|a_N|,|A|+1\},
\]
则所有 $n$ 都有 $|a_n|\le M$。这也解释了为什么必须把前面有限项一并考虑。
\textbf{常见反例：}$(-1)^n$ 有界，却不收敛：偶数子列趋于 $1$，奇数子列趋于 $-1$。而 $n$ 虽单调递增，却因无界而不收敛。
\subhead{运算与保序}
若 $a_n\to A,b_n\to B$，则和、差、积分别趋于 $A\pm B,AB$；商在 $B\ne0$ 时趋于 $A/B$。若最终 $a_n\le b_n$，则 $A\le B$。
严格不等号取极限后可能变成等号：$1/n>0$，极限却是 $0$。反过来，若 $A>0$，取误差 $A/2$，可知最终 $a_n>A/2>0$，这叫保号性。
}
% 自然分页：不人为插入分页。
% 第2页：运算与存在性，补齐典型方法。
\textcol{2}{
\subhead{有理式与根式}
\ex{例3\quad 同除最高次项}
\[
 \frac{2n^2-3n+1}{n^2+4}
 =\frac{2-3/n+1/n^2}{1+4/n^2}\longrightarrow2.
\]
分子、分母同除最高次幂；分母的极限非零，才能使用商的极限法则。
\ex{例4\quad 根式相减}
\[
 \begin{aligned}
 \sqrt{n^2+n}-n
 &=\frac{n}{\sqrt{n^2+n}+n}\\
 &=\frac1{\sqrt{1+1/n}+1}\longrightarrow\frac12.
 \end{aligned}
\]
两项都趋于无穷时，不能直接求差；先有理化，再取极限。
\subhead{夹逼：先找统一的上、下界}
\graybox{\modulehead{定理 1.7}{夹逼定理}}{
若数列 $\{b_n\}$ 和 $\{c_n\}$ 都收敛于 $a$，且对所有充分大的 $n$，有 $b_n\le a_n\le c_n$，则数列 $\{a_n\}$ 也收敛，而且极限也为 $a$。}
因为 $|\sin n|\le1$，有 $|\sin n/n|\le1/n\to0$，所以 $\sin n/n\to0$。这里不需要 $\sin n$ 自己收敛。
\ex{例5\quad 项数随 $n$ 增加的和}
设 $S_n=\sum_{k=1}^{n}(n^2+k)^{-1/2}$。对每一项都有
\[
 \frac1{\sqrt{n^2+n}}\le\frac1{\sqrt{n^2+k}}
 \le\frac1{\sqrt{n^2+1}}.
\]
相加后得到
\[
 \frac1{\sqrt{1+1/n}}\le S_n
 \le\frac1{\sqrt{1+1/n^2}},
\]
故 $S_n\to1$。不能因为“每项趋于零”就断言整个和趋于零，因为项数也在增加。
\subhead{Stolz 定理}
\lead{定理 1.21}（$\infty/\infty$ 型 Stolz 定理）设 $\{a_n\},\{b_n\}$ 是两个数列，且 $\{b_n\}$ 严格递增趋于 $+\infty$。如果
\[
 \lim_{n\to\infty}\frac{a_{n+1}-a_n}{b_{n+1}-b_n}=A,
\]
则有 $\lim_{n\to\infty}\frac{a_n}{b_n}=A$，其中 $A$ 可以是实数，也可以是 $+\infty$ 或 $-\infty$。
}
% 自然分页：不人为插入分页。
\textcol{2}{
\ex{例6\quad 用 Stolz 求极限}
令 $a_n=1+2+\cdots+n$，$b_n=n^2$。$b_n$ 严格递增且趋于无穷，又
\[
 \frac{a_{n+1}-a_n}{b_{n+1}-b_n}
 =\frac{n+1}{2n+1}\to\frac12,
\]
故 $a_n/b_n\to1/2$。此题也能用求和公式；Stolz 更适合没有现成求和公式的情况。
\subhead{单调有界与递推数列}
\graybox{\modulehead{定理 1.13}{单调有界定理}}{
单调增（减）有界数列必收敛，其极限值等于数列构成的数集的上确界（或下确界），也称为数列的上确界（或下确界）。}
\ex{例7\quad $a_1=\sqrt2,\ a_{n+1}=\sqrt{a_n+2}$}
\prooflead{证明·第一步：有界性}$0<a_1<2$；若 $0<a_n<2$，则 $0<\sqrt{a_n+2}<2$。由归纳法，所有项都在 $(0,2)$ 中。
\prooflead{第二步：单调性}$a_2>a_1$。若 $a_n>a_{n-1}$，平方根函数递增给出
\[
 a_{n+1}=\sqrt{a_n+2}>\sqrt{a_{n-1}+2}=a_n.
\]
故数列严格递增。
\prooflead{第三步：极限值}由单调有界定理，数列收敛。设极限为 $L$，由 $a_n\ge\sqrt2$ 知 $L\ge\sqrt2$。对 $a_{n+1}^2=a_n+2$ 取极限：
\[
 L^2=L+2,\qquad (L-2)(L+1)=0.
\]
排除 $L=-1$，得 $L=2$。
\graybox{易错点}{
递推式只能在已经证明收敛后用来求极限。方程的根只是候选值，还要结合正负、上下界或初值筛选。}
\subhead{根式极限与常见结论}
固定 $a>0$ 时 $\sqrt[n]{a}\to1$；还有 $\sqrt[n]{n}\to1$。后一结论可用二项式估计：设 $\sqrt[n]{n}=1+h_n$，则 $h_n\ge0$，且 $n\ge\binom n2h_n^2$，故
\[
 0\le h_n\le\sqrt{\frac2{n-1}}\longrightarrow0.
\]
另一个重要极限是 $(1+1/n)^n\to e$。一般地，
\[
 \left(1+\frac cn\right)^n\to e^c
\]
（$c$ 为常数，底数最终为正）。注意指数与括号内的小量要一起看，不能把底数先替成 $1$。
}
% 自然分页：不人为插入分页。
% 第3页：函数极限、连续及间断点。
\textcol{3}{
\subhead{函数极限与左右极限}
\graybox{\modulehead{定义 1.25}{函数极限}}{
设 $f(x)$ 在 $x_0$ 附近有定义（在 $x_0$ 不要求有定义）。如果对任意给定的正数 $\eps$，存在 $\delta>0$，当 $0<|x-x_0|<\delta$ 时有 $|f(x)-l|<\eps$。那么称 $l$ 为当 $x$ 趋向 $x_0$ 时 $f(x)$ 的极限，记成
\[
 \lim_{x\to x_0}f(x)=l,\quad\text{或 }f(x)\to l\ (x\to x_0).
\]}
$0<|x-x_0|$ 排除了点 $x_0$。所以极限只看附近的值，$f(x_0)$ 可以没有定义，也可以不等于极限。
若定义域在 $x_0$ 两侧都有点，则有限的双侧极限存在，当且仅当左、右极限都存在且相等。
\ex{例8\quad 用定义证明 $x^2\to4$（$x\to2$）}
误差是 $|x^2-4|=|x-2||x+2|$。先要求 $|x-2|<1$，便有 $|x+2|<5$。任给 $\eps>0$，取
\[
 \delta=\min\{1,\eps/5\}.
\]
当 $0<|x-2|<\delta$ 时，$|x^2-4|<5\delta\le\eps$，故结论成立。这里取最小值，是为了同时满足“控制另一个因子”和“控制误差”。
\subhead{用数列判断函数极限}
\lead{定理 1.32}函数 $f(x)$ 在 $x\to x_0$ 时有极限 $l$ 的充分必要条件是：对于任意一个以 $x_0$ 为极限的数列 $\{a_n\}$（$a_n\ne x_0$），都有 $\lim_{n\to\infty}f(a_n)=l$。
证明极限不存在时，只需找两列趋于同一点的自变量，让对应函数值趋于不同的数。例如 $\sin(1/x)$ 在 $0$ 附近可沿两列点分别取值 $1$ 和 $-1$，故极限不存在。
\subhead{两个重要极限及等价无穷小}
\[
 \frac{\sin x}{x}\to1,\qquad (1+x)^{1/x}\to e
 \quad(x\to0).
\]
角度用弧度。常用推论是
\[
 \ln(1+x)\sim x,\quad e^x-1\sim x,\quad
 1-\cos x\sim x^2/2.
\]
$\alpha\sim\beta$ 表示 $\alpha/\beta\to1$，适合在乘除中替换。例如
\[
 \frac{\sin3x}{\ln(1+2x)}
 =\frac{\sin3x}{3x}\frac{2x}{\ln(1+2x)}
   \frac32\longrightarrow\frac32.
\]
相减时不要随意替换；如 $\sin x-x$ 的两项首项抵消，需要更精细的方法。
}
% 自然分页：不人为插入分页。
\textcol{3}{
\head{2\quad 一元函数的连续性}
\subhead{连续与极限的区别}
\graybox{\modulehead{定义 2.1}{连续性}}{
设 $y=f(x)$ 在 $x_0$ 的邻域内（即包含 $x_0$ 的一个开区间）有定义，若
\[
 \lim_{x\to x_0}f(x)=f(x_0).
\]
则称 $y=f(x)$ 在 $x_0$ 处连续，$x_0$ 是 $f(x)$ 的连续点，否则称 $f(x)$ 在点 $x_0$ 处不连续，或 $x_0$ 是 $f(x)$ 的间断点。}
闭区间端点的连续性，只考察区间内的一侧。
计算连续性题，要把“附近趋于什么”和“这个点取什么值”分别写出来，再比较。多项式、指数、三角函数等在各自定义域内连续；商还要求分母非零。
\ex{例9\quad 补一个值使函数连续}
设 $f(x)=(x^2-1)/(x-1)$（$x\ne1$），并规定 $f(1)=a$。
当 $x\ne1$ 时，约分得 $f(x)=x+1$，所以 $\lim_{x\to1}f(x)=2$。连续要求 $f(1)=a=2$。约分只在 $x\ne1$ 时成立，不能据此直接代出原式在 $1$ 的值。
\ex{例10\quad 两侧公式不同}
若左侧 $f(x)=x+1$、右侧 $f(x)=x^2+1$，且 $f(0)=a$，则左右极限均为 $1$，故连续当且仅当 $a=1$。若把右侧改成 $x^2+2$，左右极限不等，怎样选 $a$ 都不能使它连续。
\subhead{间断点的分类}
\textbf{可去间断：}左右极限相同且有限，但函数值缺失或不等于极限。调整这一个点的值就能补成连续。
\textbf{跳跃间断：}左右极限均有限，但不相等。例如 $|x|/x$ 在 $0$ 的左右极限分别是 $-1,1$。
\textbf{无穷间断：}至少一侧趋于正或负无穷。例如 $1/x$ 在 $0$。
\textbf{振荡间断：}如 $\sin(1/x)$ 在 $0$，它有界，但极限不存在。不能把一切“无极限”都叫无穷间断。
\ex{例11\quad 振荡乘一个小量}
令 $f(x)=x\sin(1/x)$（$x\ne0$），$f(0)=0$。由于
\[
 |x\sin(1/x)|\le|x|\to0,
\]
故极限为 $0=f(0)$，函数在 $0$ 连续。$\sin(1/x)$ 本身振荡，但振幅被 $|x|$ 压到零。
\textbf{易错点：}“初等函数连续”限于定义域内。分段函数每段内部通常可直接判断，拼接点须另算左右极限。
}
% 自然分页：不人为插入分页。
% 第4页：定理解释、存在性证明及一致连续。
\textcol{4}{
\subhead{零点定理与介值定理}
\graybox{\modulehead{定理 2.8}{零点定理}}{
设 $f(x)\in C[a,b]$，且函数在两个端点的值 $f(a)$ 和 $f(b)$ 异号，即 $f(a)f(b)<0$，则必有一点 $\xi\in(a,b)$，使 $f(\xi)=0$。}
\lead{定理 2.9}（介值定理）设 $f(x)\in C[a,b]$，且 $f(a)\ne f(b)$，则 $f(x)$ 在 $[a,b]$ 上能取到介于 $f(a)$ 和 $f(b)$ 之间的任意值。
\prooflead{记号}$C[a,b]$ 表示定义在闭区间 $[a,b]$ 上的连续函数全体。
\ex{例12\quad 证明 $x^3+x-1=0$ 有根}
取 $F(x)=x^3+x-1$。$F$ 是多项式，在 $[0,1]$ 连续；$F(0)=-1<0$，$F(1)=1>0$。由零点定理，存在 $c\in(0,1)$ 使 $F(c)=0$。
如果还问唯一性，可对 $y>x$ 计算
\[
 F(y)-F(x)=(y-x)(y^2+xy+x^2+1)>0.
\]
所以 $F$ 严格递增，零点至多一个。
\ex{例13\quad 先构造辅助函数}
设连续函数 $f:[0,1]\to[0,1]$，证明存在 $c$ 使 $f(c)=c$。令 $g(x)=f(x)-x$，则
\[
 g(0)=f(0)\ge0,\qquad g(1)=f(1)-1\le0.
\]
若某端点等于零，直接取该端点；否则端点异号，由零点定理在 $(0,1)$ 内有零点。因而总有 $c\in[0,1]$ 满足 $f(c)=c$。
\subhead{有界性与最值定理}
\lead{定理 2.10}若 $f(x)\in C[a,b]$，则它在整个区间上有界。即存在一个常数 $M>0$，使得当 $a\le x\le b$ 时，有 $|f(x)|\le M$。
\graybox{\modulehead{定理 2.11}{最值定理}}{
若 $f(x)\in C[a,b]$，则 $f(x)$ 在 $[a,b]$ 上能取到最大值和最小值。即存在 $x_*\in[a,b]$ 和 $x^*\in[a,b]$，使得对所有的 $x\in[a,b]$，有 $f(x_*)\le f(x)\le f(x^*)$。}
\ex{例14\quad 正函数为什么能统一远离零}
若 $f$ 在 $[a,b]$ 连续且处处 $f(x)>0$，最值定理给出某个 $c$，使 $m=f(c)=\min f>0$。因此对所有 $x$，都有 $f(x)\ge m$，从而 $1/f(x)\le1/m$。
闭区间条件不能随意删：$f(x)=x$ 在 $(0,1)$ 上处处正，但没有统一的正下界；它的上确界为 $1$，却没有最大值。
}
% 自然分页：不人为插入分页。
\textcol{4}{
\subhead{一致连续}
普通连续：先固定 $x_0$，再根据 $x_0$ 和误差 $\eps$ 选 $\delta$。一致连续要求同一个 $\delta$ 对整个定义域都能用。
\graybox{\modulehead{定义 2.13}{一致连续}}{
设 $y=f(x)$ 在区间 $I$ 上有定义，对任给的 $\eps>0$，如果存在 $\delta>0$，使得任取一点 $x_0\in I$，只要 $|x-x_0|<\delta$，就有 $|f(x)-f(x_0)|<\eps$。则称函数 $f$ 在 $I$ 上是一致连续的。}
\ex{例15\quad $x^2$ 在有界区间上一致连续}
对 $x,y\in[0,M]$（$M>0$），
\[
 |x^2-y^2|=|x-y||x+y|\le2M|x-y|.
\]
任给 $\eps>0$，取 $\delta=\eps/(2M)$。当 $|x-y|<\delta$ 时，误差小于 $\eps$。这个 $\delta$ 对区间内所有点都有效。
\ex{例16\quad $x^2$ 在 $\R$ 上不一致连续}
取 $x_n=n,\ y_n=n+1/n$，则 $|x_n-y_n|=1/n\to0$，但
\[
 |x_n^2-y_n^2|=2+\frac1{n^2}>2.
\]
若一致连续，取 $\eps=1$ 所得到的 $\delta$ 应对这两列点也有效；充分大的 $n$ 满足距离小于 $\delta$，函数值之差却大于 $1$，矛盾。
\subhead{闭区间连续函数的一致连续性}
\lead{定理 2.14}有限闭区间 $[a,b]$ 上定义的连续函数 $f(x)$，一定在 $[a,b]$ 上一致连续。
\prooflead{证明思路}假设不一致连续，则存在 $\eps_0>0$ 和两列点 $x_n,y_n\in[a,b]$，满足
\[
 |x_n-y_n|<1/n,\qquad
 |f(x_n)-f(y_n)|\ge\eps_0.
\]
有界数列有收敛子列，可取 $x_{n_k}\to c\in[a,b]$。因两点距离趋零，$y_{n_k}\to c$；由连续性，两列函数值都趋于 $f(c)$，它们的差应趋零，矛盾。
\graybox{要点}{
闭区间上的连续函数：有界、取得最值、具有介值性、一致连续。做题时看清所需的是哪一个结论，写出对应条件。}
}
% 自然分页：不人为插入分页。
% 定理穿插在对应知识点中，不单独列总览。

% 习题部分从双面打印的正面页开始。
\evenbreak
% 四套题，每套一栏；伸缩留白按解答量分配。
\quizcol{5}{小测一\quad 数列极限}{
\q{1}{求 $\displaystyle\limn\frac{5n^3-n+1}{2n^3+7n^2-1}$。}
\work{2}
\q{2}{求 $\displaystyle\limn(\sqrt{n^2+3n}-n)$。}
\work{3}
\q{3}{判断 $a_n=(-1)^n$ 是否收敛，写出理由。}
\work{2}
\q{4}{用夹逼定理求 $\displaystyle\limn\frac{\sin(n^2)}{\sqrt n}$。}
\work{3}
}
% 自然分页：不人为插入分页。
\quizcol{5}{小测二\quad 极限与连续}{
\q{1}{求 $\displaystyle\lim_{x\to0}\frac{\sqrt{1+x}-1}{x}$。}
\work{2}
\q{2}{设 $f(x)=x\sin(1/x)$（$x\ne0$），$f(0)=0$。判断 $f$ 在 $0$ 处是否连续。}
\work{3}
\q{3}{判断 $f(x)=1/(x-2)$ 在 $x=2$ 的间断类型，并写出左右极限。}
\work{2}
\q{4}{设 $f(x)=(x^2-1)/(x-1)$（$x\ne1$），$f(1)=a$。求使 $f$ 在 $1$ 连续的 $a$。}
\work{3}
}
% 自然分页：不人为插入分页。
\quizcol{6}{小测三\quad 连续函数的性质}{
\q{1}{证明方程 $x^3+x-1=0$ 在 $(0,1)$ 内至少有一个实根。}
\work{4}
\q{2}{$f$ 在 $[0,1]$ 连续，且 $f(0)=-2,f(1)=3$。能否断定它在 $(0,1)$ 内有零点？说明依据。}
\work{3}
\q{3}{若 $f$ 在 $[0,1]$ 连续，能否断定它取得最大值和最小值？若区间改成 $(0,1)$，结论是否仍必然成立？举例说明。}
\work{4}
}
% 自然分页：不人为插入分页。
\quizcol{6}{小测四\quad 综合练习}{
\q{1}{设 $a_1=1,\ a_{n+1}=(a_n+2)/2$。证明数列收敛并求极限。}
\work{4}
\q{2}{设 $a_n>0$ 且 $a_n\to A>0$。证明 $\sqrt{a_n}\to\sqrt A$。}
\work{3}
\q{3}{讨论下列函数在 $0$ 处的连续性：
\[
 f(x)=\begin{cases}x+1,&x<0,\\ a,&x=0,\\ x^2+1,&x>0.\end{cases}
\]}
\work{2}
\q{4}{写出 $a_n\to A$ 的定义，并用定义证明 $1/n\to0$。}
\work{4}
}
% 自然分页：不人为插入分页。
% 答案部分从双面打印的正面页开始，给出标准过程。
\evenbreak
\textcol{7}{
\head{参考答案}
\subhead{小测一}
\ex{1.\quad 有理式极限}
分子、分母同除 $n^3$，得
\[
 \frac{5n^3-n+1}{2n^3+7n^2-1}
 =\frac{5-n^{-2}+n^{-3}}{2+7n^{-1}-n^{-3}}.
\]
当 $n\to\infty$ 时，分子趋于 $5$，分母趋于 $2\ne0$。由商的极限法则，所求极限为 $\boxed{5/2}$。
\ex{2.\quad 根式极限}
分子乘共轭，整理为
\[
 \begin{aligned}
 \sqrt{n^2+3n}-n
 &=\frac{(n^2+3n)-n^2}{\sqrt{n^2+3n}+n}\\
 &=\frac{3}{\sqrt{1+3/n}+1}.
 \end{aligned}
\]
分母趋于 $2$，所以极限为 $\boxed{3/2}$。
\ex{3.\quad 判断收敛}
对任意正整数 $k$，$a_{2k}=1$，$a_{2k-1}=-1$，故
\[
 \lim_{k\to\infty}a_{2k}=1,\qquad
 \lim_{k\to\infty}a_{2k-1}=-1.
\]
若原数列收敛，则所有子列必须收敛到同一个极限；上式矛盾。因此 $\boxed{a_n\text{ 不收敛}}$。
\ex{4.\quad 夹逼}
由 $-1\le\sin(n^2)\le1$ 且 $\sqrt n>0$，得
\[
 -\frac1{\sqrt n}\le
 \frac{\sin(n^2)}{\sqrt n}\le\frac1{\sqrt n}.
\]
两侧极限都为 $0$。由夹逼定理，所求极限为 $\boxed0$。
\prooflead{书写提醒}计算题保留关键变形；证明发散时明确指出冲突的是“所有子列应有相同极限”；夹逼题要把上下界和它们的极限一并写出。
\prooflead{指数式补充}$|q|<1$ 时 $q^n\to0$；固定 $a>1,k>0$ 时 $n^k/a^n\to0$。例如同除 $5^n$：
\[
 \frac{2^n+3\cdot4^n}{5^n+2^n}
 =\frac{(2/5)^n+3(4/5)^n}{1+(2/5)^n}\to0.
\]
\subhead{定理补充：实数的完备性}
\lead{定理 1.11}$\R$ 中任何有上（下）界的非空数集一定有上（下）确界。
\lead{定理 1.16}从任何有界的数列中可选出一个收敛的子列。
}
% 自然分页：不人为插入分页。
\textcol{7}{
\subhead{小测二}
\ex{1.\quad 有理化}
当 $x\ne0$ 且 $x$ 充分接近 $0$ 时，
\[
 \begin{aligned}
 \frac{\sqrt{1+x}-1}{x}
 &=\frac{1+x-1}{x(\sqrt{1+x}+1)}\\
 &=\frac1{\sqrt{1+x}+1}.
 \end{aligned}
\]
由初等函数连续性与商的极限法则，所求极限为 $\boxed{1/2}$。
\ex{2.\quad 用夹逼证明连续}
对 $x\ne0$，有
\[
 0\le|x\sin(1/x)|\le|x|.
\]
当 $x\to0$ 时，右侧趋于 $0$，故
\[
 \lim_{x\to0}f(x)=0.
\]
又 $f(0)=0$，极限等于函数值，因此 $\boxed{f\text{ 在 }0\text{ 连续}}$。
如需按定义写：任给 $\eps>0$，取 $\delta=\eps$，当 $|x|<\delta$ 时，$|f(x)-f(0)|\le|x|<\eps$，也得到同一结论。
\ex{3.\quad 间断点分类}
当 $x\to2^+$ 时，分母从正侧趋零；当 $x\to2^-$ 时，分母从负侧趋零。因此
\[
 \lim_{x\to2^+}\frac1{x-2}=+\infty,\qquad
 \lim_{x\to2^-}\frac1{x-2}=-\infty.
\]
所以 $\boxed{x=2\text{ 是无穷间断点}}$。
仅写“分母为零”还不足以分类，因为分母为零也可能对应可去间断。
\ex{4.\quad 补点连续}
$x\ne1$ 时，
\[
 f(x)=\frac{(x-1)(x+1)}{x-1}=x+1.
\]
故 $\lim_{x\to1}f(x)=2$。连续的充要条件为 $f(1)=a=2$，所以 $\boxed{a=2}$。
\graybox{易错点}{
约去 $x-1$ 只是在去心邻域内化简，并没有自动给原函数在 $1$ 处定义数值。最后仍要单独检查 $f(1)$。}
\subhead{自查}
有界不一定收敛；连续不一定一致连续；开区间上的连续函数未必取得最值。各举一个反例，并指出不能使用相应定理的原因。
}
% 自然分页：不人为插入分页。
\textcol{8}{
\subhead{小测三}
\ex{1.\quad 零点存在性}
令 $F(x)=x^3+x-1$。由于 $F$ 为多项式，它在 $[0,1]$ 上连续，且
\[
 F(0)=-1<0,\qquad F(1)=1>0.
\]
由零点定理，存在 $c\in(0,1)$ 使 $F(c)=0$，即 $c^3+c-1=0$。故方程在该开区间内至少有一个实根。
题目没有要求求出根，也没有要求证明唯一性；以上过程已完成所问。
\ex{2.\quad 抽象函数的零点}
由题设，$f$ 在闭区间 $[0,1]$ 连续，且
\[
 f(0)f(1)=(-2)\cdot3=-6<0.
\]
由零点定理，存在 $c\in(0,1)$，使 $f(c)=0$。所以答案是“能”。这里不需要知道 $f$ 的具体表达式。
\ex{3.\quad 最值与条件}
因为 $f$ 在闭有界区间 $[0,1]$ 连续，由最值定理，存在 $u,v\in[0,1]$，使对任意 $x\in[0,1]$，
\[
 f(u)\le f(x)\le f(v).
\]
其中 $f(u)$ 是最小值，$f(v)$ 是最大值。
若改为 $(0,1)$，结论不再必然成立。例如 $f(x)=x$ 在 $(0,1)$ 连续，但对任意 $x$，
\[
 0<x/2<x<(x+1)/2<1.
\]
区间里总有取值更小和更大的点，因此既没有最小值，也没有最大值。
\graybox{条件检查}{
零点定理：闭区间连续，端点严格异号。最值定理：闭有界区间上连续。缺少条件时，不能直接套定理；这不等于每个具体函数一定不满足结论。}
\subhead{定理补充：基本列}
\lead{定义 1.17}数列 $\{a_n\}$ 称为基本列，若对任意给定的正数 $\eps$，存在整数 $N=N(\eps)$（即 $N$ 可能依赖于 $\eps$），使得当 $m,n>N$ 时，就有 $|a_n-a_m|<\eps$。
\lead{定理 1.18}（Cauchy 收敛准则）数列 $\{a_n\}$ 收敛的充分必要条件是 $\{a_n\}$ 是基本列。
}
% 自然分页：不人为插入分页。
\textcol{8}{
\subhead{小测四}
\ex{1.\quad 递推数列}
先证 $0<a_n<2$。$a_1=1$ 满足；若 $0<a_n<2$，则 $1<a_{n+1}=(a_n+2)/2<2$。由归纳法，结论对所有 $n$ 成立。于是
\[
 a_{n+1}-a_n=\frac{2-a_n}{2}>0.
\]
数列递增且有上界 $2$，故收敛。设极限为 $L$，对递推式取极限得 $L=(L+2)/2$，故 $\boxed{L=2}$。
另解可用 $2-a_{n+1}=(2-a_n)/2$，推出 $2-a_n=2^{1-n}\to0$；这也直接证明了收敛性。
\ex{2.\quad 平方根的极限}
由 $a_n>0,A>0$，有

undefined
undefined
undefined
undefined
undefined
\[
 \begin{aligned}
 |\sqrt{a_n}-\sqrt A|
 &=\frac{|a_n-A|}{\sqrt{a_n}+\sqrt A}\\
 &\le\frac{|a_n-A|}{\sqrt A}\longrightarrow0.
 \end{aligned}
\]
由夹逼定理，$\boxed{\sqrt{a_n}\to\sqrt A}$。这个估计不预先使用待证的平方根极限，因而没有循环论证。
\ex{3.\quad 分段函数连续性}
分别计算：
\[
 \lim_{x\to0^-}(x+1)=1,\qquad
 \lim_{x\to0^+}(x^2+1)=1.
\]
左右极限相等，所以 $\lim_{x\to0}f(x)=1$。又 $f(0)=a$，故连续当且仅当 $\boxed{a=1}$；$a\ne1$ 时为可去间断点。
\ex{4.\quad 定义证明}
$a_n\to A$ 的定义为：对任意 $\eps>0$，存在正整数 $N$，使 $n>N$ 时 $|a_n-A|<\eps$。
对 $a_n=1/n$，任给 $\eps>0$，取
\[
 N=\lfloor1/\eps\rfloor+1.
\]
若 $n>N$，则 $n>1/\eps$，所以
\[
 \left|\frac1n-0\right|=\frac1n<\eps.
\]
由定义，$\boxed{\limn1/n=0}$。$N$ 是根据给定误差选出的统一项号；写出“任给 $\eps$、取 $N$、当 $n>N$”三个环节即可。
\prooflead{基本列的理解}基本列也称柯西列，判断时不用预先知道极限。仅有 $a_{n+1}-a_n\to0$ 不够，例如 $a_n=\sqrt n$ 仍趋于无穷。
}
\evenfinish
\end{document}

